% Abstracts in English and Polish. Write them last.
% Outline: a one-page summary: the problem, the approach (an SSVEP speller on a low-cost
% dry-electrode cap, supported by a language model) and the main results.
% Don't use \ac here: the abstracts define their own abbreviations in plain text.
\begin{abstracts}

  \AbstractENG

  This thesis concerns the design and implementation of a non-invasive, real-time
  brain–computer interface (BCI) intended to support communication for people with severe
  motor impairments. The system is based on the steady-state visual evoked potential
  (SSVEP) paradigm. The user selects characters or words by focusing their gaze on
  interface elements that flicker at different frequencies.

  The engineering part includes several components: a graphical user interface with
  precise visual stimulus timing and a signal processing pipeline.

  The research part compares spatial filtering and classification methods, including
  canonical correlation analysis (CCA) and task-related component analysis (TRCA),
  together with selected machine learning algorithms. The algorithms were trained and
  evaluated on offline recorded data. The system is integrated with a natural language
  processing (NLP) model that predicts subsequent words in real time. This is intended to
  increase communication speed and reduce user fatigue.

  The integrated system was validated in online experiments. Classification accuracy and
  information transfer rate (ITR) were measured both with and without text prediction. The
  results show to what extent combining SSVEP classification with language prediction
  improves communication efficiency in BCI systems.

  \AbstractPL

  Praca dotyczy projektu i implementacji nieinwazyjnego interfejsu mózg–komputer (BCI)
  działającego w czasie rzeczywistym, przeznaczonego do komunikacji osób z ciężkimi
  dysfunkcjami ruchowymi. System wykorzystuje paradygmat potencjałów wywołanych stanu
  ustalonego (SSVEP). Użytkownik wybiera w nim znaki lub słowa, skupiając wzrok na
  elementach interfejsu migających z różnymi częstotliwościami.

  Część inżynierska obejmuje kilka elementów: graficzny interfejs użytkownika z precyzyjną
  synchronizacją bodźców wzrokowych oraz potok przetwarzania sygnału.

  Część badawcza polega na porównaniu metod filtracji przestrzennej i klasyfikacji, m.in.
  analizy korelacji kanonicznej (CCA) i analizy komponentów powiązanych z zadaniem (TRCA),
  a także wybranych algorytmów uczenia maszynowego. Algorytmy wytrenowano i oceniono na
  danych zarejestrowanych offline. System zintegrowano z modelem przetwarzania języka
  naturalnego (NLP), który w czasie rzeczywistym podpowiada kolejne słowa. Ma to zwiększyć
  szybkość komunikacji i ograniczyć zmęczenie użytkownika.

  Działanie zintegrowanego systemu zweryfikowano w eksperymentach online. Oceniono
  dokładność klasyfikacji oraz szybkość transferu informacji (ITR) zarówno z predykcją
  tekstu, jak i bez niej. Wyniki pozwalają określić, w jakim stopniu połączenie
  klasyfikacji sygnałów SSVEP z predykcją językową poprawia efektywność komunikacji w
  systemach BCI.

\end{abstracts}